\documentclass[12pt,a4paper]{article}

\usepackage[utf8]{inputenc}
\usepackage[T1]{fontenc}
\usepackage[provide=*,spanish]{babel}
\usepackage{mathpazo}
\usepackage{microtype}
\usepackage{amsmath,amssymb,amsthm}
\usepackage{geometry}
\usepackage{enumitem}
\usepackage[hidelinks]{hyperref}

\geometry{margin=2.5cm}

\title{Ecuaciones Diferenciales Ordinarias\\Clase V -- Teórico}
\author{Benjamín Marcolongo}
\date{}

\begin{document}

	\maketitle

	\section*{Introducción}

	En la Clase IV estudiamos las ecuaciones lineales de primer orden y el método del factor integrante. En esta clase utilizaremos esas herramientas para resolver una familia de ecuaciones no lineales: las ecuaciones de Bernoulli.

	La idea central no consiste en resolver directamente la ecuación no lineal, sino en introducir una nueva variable dependiente que satisfaga una ecuación lineal. Por eso, el método de Bernoulli es una reducción a un problema ya conocido.

	En todo el material, \(x\) representa la variable independiente e \(y=y(x)\) la función desconocida. Escribiremos \(dy/dx\) para la derivada de \(y\) respecto de \(x\).

	\clearpage
	\section{Ecuaciones de Bernoulli}

	Una ecuación de \textbf{Bernoulli} tiene la forma

	\[
	\boxed{
	\frac{dy}{dx}+p(x)y=q(x)y^n,
	}
	\]
	%
	donde \(p\) y \(q\) son funciones conocidas de \(x\) y \(n\) es una constante real. Si \(n=0\), la ecuación ya es lineal. Si \(n=1\), también es lineal y además es separable. El caso nuevo corresponde a \(n\neq0\) y \(n\neq1\).

	Antes de realizar cualquier división por una potencia de \(y\), debe verificarse por separado si \(y=0\) satisface la ecuación original. La respuesta depende del valor de \(n\) y del dominio donde esté definida la potencia \(y^n\).

	\subsection{Reducción a una ecuación lineal}

	En un intervalo donde la solución no se anula y las potencias involucradas están definidas, dividimos por \(y^n\):

	\[
	y^{-n}\frac{dy}{dx}+p(x)y^{1-n}=q(x).
	\]
	%
	Definimos la nueva variable dependiente

	\[
	u(x)=y(x)^{1-n}.
	\]
	%
	Por la regla de la cadena,

	\[
	\frac{du}{dx}
	=
	(1-n)y^{-n}\frac{dy}{dx}.
	\]
	%
	Multiplicando la ecuación anterior por \(1-n\), obtenemos

	\[
	\boxed{
	\frac{du}{dx}
	+(1-n)p(x)u
	=
	(1-n)q(x),
	}
	\]
	%
	que es una ecuación lineal de primer orden para la función desconocida \(u=u(x)\). Se resuelve con un factor integrante y, finalmente, se vuelve a la variable \(y\).

	\subsection{Procedimiento}

	\begin{enumerate}[label=\arabic*.]
		\item escribir la ecuación en la forma de Bernoulli;
		\item identificar \(p(x)\), \(q(x)\) y el exponente \(n\);
		\item comprobar directamente las soluciones que puedan perderse al dividir por \(y^n\);
		\item introducir \(u=y^{1-n}\);
		\item resolver la ecuación lineal obtenida para \(u\);
		\item regresar a la variable \(y\);
		\item imponer la condición inicial y determinar el intervalo de definición.
	\end{enumerate}

	\clearpage
	\subsection{Ejemplo}

	Consideremos

	\[
	t^2\frac{dy}{dt}+y^2-ty=0.
	\]
	%
	Ahora \(t\) es la variable independiente e \(y=y(t)\) la función desconocida. Como el coeficiente principal se anula en \(t=0\), trabajaremos en un intervalo que no contenga ese punto. Dividiendo por \(t^2\),

	\[
	\frac{dy}{dt}-\frac{1}{t}y
	=
	-\frac{1}{t^2}y^2.
	\]
	%
	Esta es una ecuación de Bernoulli con

	\[
	p(t)=-\frac{1}{t},
	\qquad
	q(t)=-\frac{1}{t^2},
	\qquad
	n=2.
	\]
	%
	La función \(y(t)=0\) es una solución y debe conservarse. Para las soluciones no nulas, definimos

	\[
	u(t)=y(t)^{-1}=\frac{1}{y(t)}.
	\]
	%
	Como \(1-n=-1\), la ecuación transformada es

	\[
	\frac{du}{dt}+\frac{1}{t}u=\frac{1}{t^2}.
	\]
	%
	En el intervalo \(t>0\), el factor integrante es

	\[
	\mu(t)=\exp\left(\int\frac{1}{t}\,dt\right)=t.
	\]
	%
	Por lo tanto,

	\[
	\frac{d}{dt}\bigl(tu(t)\bigr)=\frac{1}{t}.
	\]
	%
	Integrando,

	\[
	tu(t)=\ln t+C,
	\]
	%
	donde \(C\) es una constante real. Como \(u=1/y\), obtenemos

	\[
	\boxed{
	y(t)=\frac{t}{\ln t+C},
	\qquad t>0.
	}
	\]
	%
	Esta expresión es válida en cualquier intervalo contenido en \(t>0\) donde \(\ln t+C\neq0\). En un intervalo contenido en \(t<0\), la expresión correspondiente utiliza \(\ln|t|\). La familia de soluciones no nulas no incluye \(y=0\), que debe agregarse por separado.

	\section{Cómo reconocer el método}

	Una ecuación de Bernoulli combina un término derivativo, un término lineal en la función desconocida y una potencia de esa función:

	\[
	\frac{dy}{dx}+p(x)y=q(x)y^n.
	\]
	%
	Antes de aplicar el cambio de variable conviene verificar si la ecuación también es separable o si los valores \(n=0\) o \(n=1\) la reducen inmediatamente a un caso más simple.

	\begin{quote}
		\textbf{Idea central.} El cambio \(u=y^{1-n}\) no elimina la no linealidad por casualidad. Su derivada contiene precisamente el producto \(y^{-n}dy/dx\) que aparece al dividir la ecuación por \(y^n\).
	\end{quote}

	\section*{Articulación con la Guía 1}

	El método desarrollado permite resolver el problema 24 de la Guía 1.

	\section*{Bibliografía y recursos recomendados}

	\begin{itemize}[leftmargin=*]
		\item D. G. Zill, \textit{A First Course in Differential Equations with Modeling Applications}, 12.\textsuperscript{a} ed., sección 2.5.
	\end{itemize}

	\subsection*{Nota sobre la elaboración del material}

	Este material fue elaborado por Benjamín Marcolongo con la asistencia de \textbf{GPT-5.6 Sol}, un modelo de OpenAI utilizado a través del entorno Codex, para la organización, redacción y revisión del contenido.

\end{document}
